% Additional explanations for the middle group of lessons.

\expandafter\def\csname sourcewalkthrough@26\endcsname{%
A useful way to begin is to ask: what does the current sum
mean at one exact index? It is not the best sum anywhere in the array. It is
the best sum of a subarray that must end at the current value. This difference
is what makes one pass possible. If the old running sum is negative, carrying
it forward only makes the next choice worse, so the scan starts again. If it
is positive, adding the new value may continue a useful block.

For \texttt{[-2,1,-3,4,-1,2,1,-5,4]}, the running value is first \(-2\).
At 1, starting again is better than extending \(-2\), so the running value
becomes 1. The same decision happens again at 4 because the sum before it is
negative. From 4 onward, the values \(-1,2,1\) still leave a positive running
sum, so that block is kept. A separate variable stores the best running value
ever seen. Initializing both variables from the first element, rather than
zero, is important when every array value is negative.}

\expandafter\def\csname sourcewalkthrough@27\endcsname{%
Products need one more piece of state than sums. A very small negative product
can become a very large positive product after another negative number. For
that reason, the scan keeps both the largest and smallest product ending at
the current position. Neither one can be discarded early.

At every value there are three choices: start a new subarray with this value,
multiply it by the previous maximum, or multiply it by the previous minimum.
The largest candidate becomes the new maximum and the smallest becomes the
new minimum. With \texttt{[-2,3,-4]}, the states after 3 are maximum 3 and
minimum \(-6\). When \(-4\) arrives, \((-6)(-4)=24\), so the old minimum
creates the new maximum. Zero also works naturally: it can end the old block
and start a new one. Both old states must be saved before either is updated,
otherwise the second calculation may use a value from the current step.}

\expandafter\def\csname sourcewalkthrough@28\endcsname{%
The window does not need every character to match already. It only needs to
be repairable with at most \(k\) replacements. Inside a window of length
\(L\), keep the character that appears most often and replace the other
\(L-f\) positions. The window is allowed exactly when \(L-f\le k\).

As the right edge moves, update the count for the new character and the
largest frequency seen. If the replacement cost is too large, move the left
edge and remove those characters from the active counts. The stored maximum
frequency is allowed to be slightly old. It may delay a shrink, but it cannot
produce a better answer that was never possible earlier. For
\texttt{AABABBA} with \(k=1\), \texttt{AABA} is valid because three of its
four positions are A. Extending farther breaks the budget, so the left edge
moves. Each edge moves only forward, which is why the work is linear.}

\expandafter\def\csname sourcewalkthrough@29\endcsname{%
A palindrome grows from its middle. An odd-length palindrome has one middle
character, while an even-length palindrome has a gap between two characters.
Those are the only two kinds of centers, so testing both at every position
finds every possible palindrome.

For each center, place one pointer on the left and one on the right. While the
indexes are valid and the characters match, record the current length and
move both pointers outward. In \texttt{babad}, the center at the first
\texttt{a} expands to \texttt{bab}; the next center can produce
\texttt{aba}. Either length-three answer is valid. The important boundary
detail is that the pointers move one step beyond the palindrome before the
loop stops, so the found substring begins at \texttt{left+1} and has length
\texttt{right-left-1}. Empty input and a two-character even palindrome are
good checks for the index arithmetic.}

\expandafter\def\csname sourcewalkthrough@30\endcsname{%
The active substring must contain no repeated character. A direct solution
can restart from every index, but it repeatedly examines the same characters.
The sliding window keeps the longest useful suffix instead. A table records
the most recent position of each character.

When the right pointer sees a character that already appears inside the
window, the left pointer jumps to one position after that old occurrence. It
must never move backward, so the update uses the larger of the current left
edge and the stored position plus one. In \texttt{abba}, seeing the second
\texttt{b} moves left past the first b. When the final a is read, its old
position lies outside the current window, so left stays where it is. This is
the case that breaks solutions that assign \texttt{left=last[a]+1} without a
maximum. The answer is updated after the window is valid.}

\expandafter\def\csname sourcewalkthrough@31\endcsname{%
This problem asks for a count rather than one longest answer. Every single
character is already a palindrome, and larger palindromes are formed by
matching equal ends around a smaller palindrome. A two-dimensional table can
store whether each substring \([i,j]\) is palindromic.

The ends match when \texttt{s[i]==s[j]}. The inside is automatically valid
for substrings of length one or two; for longer substrings, the table entry
\([i+1,j-1]\) must already be true. That dependency tells us to fill shorter
lengths before longer ones. Each true entry adds one to the answer. For
\texttt{aaa}, the three single letters, two length-two substrings, and the
whole string give six. Expanding around centers is another valid approach and
uses constant extra space, but the table makes the exact recurrence visible
and connects naturally to later interval dynamic-programming problems.}

\expandafter\def\csname sourcewalkthrough@32\endcsname{%
Only one question matters after sorting: does the next meeting begin before
the current one ends? If it does, one room cannot hold both meetings. If it
does not, the current room becomes free in time.

Sorting by start time puts the meetings in the order in which they must be
considered. For \texttt{[[0,30],[5,10],[15,20]]}, the meeting beginning at 5
overlaps the one ending at 30, so the answer is false. For
\texttt{[[7,10],[2,4]]}, sorting gives \([2,4]\) then \([7,10]\), and there
is no overlap. Touching endpoints such as \([1,3]\) and \([3,5]\) are
allowed because one meeting ends exactly when the next begins. The sort costs
\(O(n\log n)\); after that, the scan only compares neighboring intervals.}

\expandafter\def\csname sourcewalkthrough@33\endcsname{%
Treat each position as a statement about future reach. From the
current index, \texttt{i+nums[i]} is the farthest place reachable with one
more jump. The algorithm keeps the farthest position reachable from any index
seen so far.

Before using an index, first check that it is not beyond the current reach. If
\texttt{i>farthest}, there is a gap that no earlier jump can cross, so the
answer is false. Otherwise, enlarge \texttt{farthest} with the current jump.
For \texttt{[2,3,1,1,4]}, index 0 reaches 2, and index 1 extends the reach to
4, which is the final index. For \texttt{[3,2,1,0,4]}, the reach stops at 3;
index 4 cannot be processed. The algorithm does not need to construct the
actual jump sequence because the question asks only whether a path exists.}

\expandafter\def\csname sourcewalkthrough@34\endcsname{%
The depth of a node is one plus the larger depth of its two children. This
definition leads directly to a recursive solution. A null child contributes
zero, and a leaf therefore returns one.

The calls first move down until they reach null pointers. As recursion returns,
each node receives the completed depths of its left and right subtrees and
adds one for itself. In a tree whose root has a shallow left child and a
deeper right chain, the maximum follows the right side even though both sides
are still evaluated. Every node is visited once, so the work is linear. The
extra space comes from the call stack and equals the tree height: logarithmic
for a balanced tree and linear for a completely skewed tree. A breadth-first
solution can also count levels, but recursion matches the mathematical
definition with less bookkeeping.}

\expandafter\def\csname sourcewalkthrough@35\endcsname{%
Inorder traversal of a binary search tree visits values in sorted order. The
\(k\)-th value produced by that traversal is therefore the \(k\)-th smallest
value in the tree. An explicit stack lets the method stop as soon as that
value is reached.

Starting at the root, push every node on the path to the leftmost child. Pop
one node, count it, and then repeat the same left-path process from its right
child. The stack contains ancestors whose own visit or right subtree is still
waiting. Decrement \(k\) on each pop, not on each push. When it reaches zero,
return that node's value. In a balanced tree the stack holds at most the tree
height. The worst case is a skewed tree, but early stopping often avoids
visiting every node.}

\expandafter\def\csname sourcewalkthrough@36\endcsname{%
The search can use the ordering of the BST rather than building parent links.
At each node, compare both target values with the current value. If both are
smaller, both targets must lie in the left subtree. If both are larger, both
must lie in the right subtree.

The first time the targets do not share one strict direction, the paths have
split. The current node is then their lowest common ancestor. Equality also
stops the search because a node is allowed to be an ancestor of itself. For a
tree rooted at 6 with targets 2 and 4, both values first send the search left.
At node 2, one target equals the current value and the other lies to the right,
so node 2 is the answer. Only one root-to-node path is followed, giving
\(O(h)\) time and constant extra space.}

\expandafter\def\csname sourcewalkthrough@37\endcsname{%
Meeting Rooms II asks for the largest number of meetings active at the same
time. Sort the meetings by start time and keep the end time of each occupied
room in a min heap. The smallest end is the room that becomes available first.

Before placing a meeting, remove every end time that is less than or equal to
its start. Those rooms are free. Push the new end, then record the largest heap
size. With \texttt{[[0,30],[5,10],[15,20]]}, the heap becomes \([30]\), then
contains 10 and 30, so two rooms are needed. At start 15, the end 10 is
removed before 20 is inserted. Keeping only the earliest end at the top is
what makes each reuse decision efficient. A max heap would expose the wrong
room. Sorting and heap operations give \(O(n\log n)\) time.}

\expandafter\def\csname sourcewalkthrough@38\endcsname{%
The greedy choice is to keep the interval that ends first. A shorter finishing
time leaves at least as much room for every later interval. Sort by end, keep
the first interval, and compare every next start with the last kept end.

If the next interval begins too early, count one removal and keep the earlier
end already chosen. If it begins at or after that end, it is compatible and
becomes the new last interval. For \texttt{[[1,2],[2,3],[3,4],[1,3]]}, the
first three short intervals can all remain, while \([1,3]\) must be removed.
An exchange argument explains the choice: whenever another solution keeps a
later-finishing interval, replacing it with the greedy interval cannot make a
future compatible interval invalid. The scan therefore keeps as many
intervals as possible, and removals equal total intervals minus kept ones.}

\expandafter\def\csname sourcewalkthrough@39\endcsname{%
Sorting by start time turns many separate overlap checks into one comparison
with the active merged interval. The output is built from left to right. Its
last interval is the only one that can still grow.

If the next start is greater than the last output end, there is a gap, so the
next interval is appended. Otherwise the intervals overlap and the last end
becomes the larger of the two ends. For \texttt{[[1,3],[2,6],[8,10],[15,18]]},
the first two ranges become \([1,6]\). The next start, 8, lies beyond 6, so a
new block begins. Contained intervals need no special case: taking the maximum
end keeps the larger range. Equal endpoints merge because they share a point.
After sorting, no interval earlier than the last output block can overlap a
future interval.}

\expandafter\def\csname sourcewalkthrough@40\endcsname{%
Because the original intervals are already sorted and non-overlapping, they
must fall into three consecutive regions around the inserted interval. First
come intervals completely before it, then intervals touching or overlapping
it, and finally intervals completely after it.

Copy the first region unchanged. In the overlap region, repeatedly lower the
new start and raise the new end. When that loop stops, append the fully merged
new interval once, then copy the remaining suffix. Inserting \([4,8]\) into
\texttt{[[1,2],[3,5],[6,7],[9,12]]} first copies \([1,2]\), then expands the
new range to \([3,8]\), and finally copies \([9,12]\). The comparisons are
deliberately strict before the range and non-strict during merging so touching
endpoints are handled as overlaps. Each original interval is examined once.}

\expandafter\def\csname sourcewalkthrough@41\endcsname{%
The outer grid scan finds the first cell of each island. Once land is found,
DFS marks every land cell connected to it in the four allowed directions.
Changing visited land to water lets the input grid serve as the visited set.

The DFS stops at a boundary, at water, or at a cell that was already erased.
Only the outer scan increases the island count; recursive calls only finish
the current component. This separation prevents one island from being counted
many times. A diagonal land cell is not reached because the recursive calls
use only up, down, left, and right. Every cell is checked a constant number of
times, so time is \(O(mn)\). The recursion stack can hold \(O(mn)\) cells for
one large winding island; an explicit stack can be used when recursion depth
is a concern.}

\expandafter\def\csname sourcewalkthrough@42\endcsname{%
The graph may arrive as an edge list, but DFS needs to know the neighbors of
one vertex quickly. First build an adjacency list by adding each undirected
edge in both directions. Then scan all vertex numbers.

When an unvisited vertex appears, it starts one new connected component.
Increase the answer and run DFS or BFS to mark every vertex reachable from it.
Later vertices from that same component will already be marked and will not
increase the count. Isolated vertices still count because they appear in the
outer scan even though their neighbor list is empty. With \(n\) vertices and
\(e\) edges, adjacency construction and traversal take \(O(n+e)\) time. The
visited array prevents cycles from causing repeated recursion.}

\expandafter\def\csname sourcewalkthrough@43\endcsname{%
The hash set gives fast membership tests, but starting a count from every value
would still repeat work. A value begins a consecutive sequence only when its
predecessor is absent. That one test identifies the left edge of each run.

For every start, test \(x+1,x+2\), and so on until a number is missing. Record
the longest length. In \texttt{[100,4,200,1,3,2]}, values 2, 3, and 4 do not
start scans because their predecessors exist. Only 1 grows through the run
\(1,2,3,4\). Although there is a loop inside another loop, each number belongs
to only one forward scan, so the total expected work is linear. Duplicates
disappear naturally in the set. Sorting would also work, but it costs
\(O(n\log n)\) and changes the recognition pattern.}

\expandafter\def\csname sourcewalkthrough@44\endcsname{%
The dummy node removes the special case for the first output node. A tail
pointer always marks the final node in the merged list. Compare the two current
values, connect the smaller node after the tail, and advance only the list
that supplied it.

No new data nodes are required; the method reuses the existing links. When one
list ends, its competitor is already sorted, so the complete remaining suffix
can be attached in one step. For \texttt{1->2->4} and \texttt{1->3->4}, the
tail follows values 1, 1, 2, 3, 4, 4. Saving the next choice before changing
links is useful when drawing the process. Time is linear in the total number
of nodes, and apart from the dummy node the algorithm uses constant extra
space.}

\expandafter\def\csname sourcewalkthrough@45\endcsname{%
At any moment, the next output node must be the smallest current head among
all lists. A min heap stores exactly those candidates. Push every non-null
head, pop the smallest node, attach it to the result, and push that node's next
pointer if it exists.

The heap never needs every node at once; it holds at most one candidate from
each list. With \(k\) lists, each push and pop costs \(O(\log k)\). Across
\(N\) total nodes, the running time is \(O(N\log k)\), and heap space is
\(O(k)\). The comparator must order node values in the min-heap direction.
As in merging two lists, a dummy head simplifies construction and the
original nodes can be relinked instead of copied. Empty lists are skipped when
the heap is initialized.}

\expandafter\def\csname sourcewalkthrough@46\endcsname{%
The table state \(dp[i][j]\) means the LCS length between suffixes beginning at
positions \(i\) and \(j\). If the two current characters match, that character
can be used and the answer is one plus the diagonal state. If they do not
match, one of the two current characters must be skipped.

This gives the recurrence
\(dp[i][j]=1+dp[i+1][j+1]\) for a match, otherwise the maximum of
\(dp[i+1][j]\) and \(dp[i][j+1]\). Fill the table from the bottom right so
all three needed states already exist. For \texttt{abcde} and \texttt{ace},
the matching choices a, c, and e produce length three without needing to be
adjacent. Confusing a subsequence with a substring is the main conceptual
mistake. The full table costs \(O(mn)\) time and space, and rows can be rolled
when only the length is required.}

\expandafter\def\csname sourcewalkthrough@47\endcsname{%
The window is valid when it contains every required occurrence from the target,
including duplicates. A count table begins with the target requirements, and
\texttt{missing} stores how many required character occurrences are still
absent.

When the right edge adds a character whose remaining count is positive,
\texttt{missing} decreases. Every added character then lowers its count; extra
copies make the count negative. Once \texttt{missing} reaches zero, move the
left edge while the window stays valid, recording shorter answers. Removing a
character raises its count, and if it becomes positive, a required occurrence
has been lost. For \texttt{ADOBECODEBANC} and \texttt{ABC}, this process first
finds \texttt{ADOBEC} and later contracts to \texttt{BANC}. Both pointers move
only forward, so the scan is linear.}

\expandafter\def\csname sourcewalkthrough@48\endcsname{%
The expected values are \(0,1,\ldots,n\), while the array contains every one
except the missing value. XOR combines the expected values with the actual
values. Every number that appears in both groups cancels because
\(x\mathbin{\char94}x=0\), and zero does not change another value.

Start the answer at \(n\), because loop indexes supply only 0 through
\(n-1\). At each index, XOR both the index and the array value. For
\texttt{[3,0,1]}, the combined expression contains 0, 1, and 3 twice, leaving
only 2. XOR is associative and commutative, so the order does not matter.
This avoids the overflow risk of a sum formula and uses constant extra space.
The important tests are a missing zero and a missing \(n\), since they expose
incorrect initialization.}

\expandafter\def\csname sourcewalkthrough@49\endcsname{%
Subtracting one from a nonzero binary number changes its lowest 1 bit to 0 and
changes all lower 0 bits to 1. ANDing that result with the original number
therefore clears exactly the lowest set bit and leaves every higher bit alone.

Repeat \texttt{n \&= n-1} and count the iterations. Binary
\texttt{1011000} becomes \texttt{1010000}, then \texttt{1000000}, then zero,
so it contains three 1 bits. The loop runs once per set bit rather than once
per bit position. For a fixed 32-bit value, both descriptions are constant
time, but the bit-clearing form shows the useful operation directly. The input
must be unsigned so shifts or arithmetic around the high bit do not introduce
signed-number behavior. Zero correctly skips the loop and returns zero.}